\section{ابرتفکیک کور دنیای واقعی با داده‌های کاملاً مصنوعی}
\label{sec:real_esrgan}

روش‌های ابرتفکیک تصویر مبتنی بر شبکه‌های عصبی عمیق، با وجود پیشرفت‌های قابل‌توجه، عمدتاً بر فرض ساده‌سازی‌شدهٔ هستهٔ کاهش‌نمونه‌برداری دومکعبی استوارند \cite{lim2017enhanceddeepresidualnetworks, zhang2018imagesuperresolutionusingdeep, wang2018esrganenhancedsuperresolutiongenerative}. این فرض با تخریب‌های واقعی تصاویر دنیای حقیقی تفاوت فاحشی دارد و موجب ناکارآمدی این روش‌ها در سناریوهای عملی می‌شود. \lr{Wang et al.} \cite{wang2021realesrgan} با ارائهٔ \lr{Real-ESRGAN}، مدل قدرتمند \lr{ESRGAN} \cite{wang2018esrganenhancedsuperresolutiongenerative} را به یک ابزار کاربردی برای بازسازی تصاویر دنیای واقعی گسترش دادند. ایدهٔ محوری این کار، آموزش شبکه با جفت‌داده‌های کاملاً مصنوعی است که از یک فرآیند مدل‌سازی تخریب مرتبهٔ بالا تولید می‌شوند. در ادامه جزئیات این رویکرد به تفصیل مرور می‌شود.
\begin{figure}
	\centering
	\includegraphics[width=1\linewidth]{figures/ESR_1.png}
	\caption{از چپ به راست ردیف ها حاصل روش دومکعبی، ESRGAN, RealSR, Real-ESRGAN هستند.}
	\label{fig:ESR_1}
\end{figure}

\subsection{مدل تخریب کلاسیک}
\label{subsec:classical_deg}

مدل تخریب کلاسیک پایهٔ بسیاری از روش‌های ابرتفکیک کور است. در این مدل، تصویر با وضوح پایین $\mathbf{x}$ از تصویر حقیقی $\mathbf{y}$ به‌صورت زیر تولید می‌شود:
\begin{equation}
	\label{eq:classical_deg}
	\mathbf{x} = \mathcal{D}(\mathbf{y}) = \bigl[(\mathbf{y} \otimes \mathbf{k}) \!\downarrow_{r} + \mathbf{n}\bigr]_{\text{JPEG}},
\end{equation}
که $\mathbf{k}$ هستهٔ تاری، $\downarrow_{r}$ عملگر کاهش‌نمونه‌برداری با ضریب $r$، $\mathbf{n}$ نویز افزوده و $[\cdot]_{\text{JPEG}}$ فشرده‌سازی JPEG است. این مدل چهار مؤلفهٔ اصلی تخریب را شامل می‌شود که در ادامه هر یک توصیف می‌گردد.

\subsection{تاری}
تخریب تاری معمولاً به‌صورت پیچش با یک فیلتر خطی مدل می‌شود. فیلترهای گاوسی همسانگرد و ناهمسانگرد رایج‌ترین انتخاب‌ها هستند. برای هستهٔ تاری گاوسی $\mathbf{k}$ با اندازهٔ $(2t+1)\times(2t+1)$، المان $(i,j) \in [-t,t]$ به‌صورت زیر از توزیع گاوسی نمونه‌برداری می‌شود:
\begin{equation}
	\label{eq:gauss_kernel}
	k(i,j) = \frac{1}{N}\exp\!\Bigl(-\frac{1}{2}\,\mathbf{C}^{\top}\boldsymbol{\Sigma}^{-1}\mathbf{C}\Bigr), \quad \mathbf{C}=[i,\,j]^{\top},
\end{equation}
که $\boldsymbol{\Sigma}$ ماتریس کوواریانس و $N$ ثابت نرمال‌سازی است. ماتریس کوواریانس به‌صورت
\begin{equation}
	\boldsymbol{\Sigma} = \mathbf{R}
	\begin{bmatrix}
		\sigma_{1}^{2} & 0 \\
		0 & \sigma_{2}^{2}
	\end{bmatrix}
	\mathbf{R}^{\top},
\end{equation}
تجزیه می‌شود که $\mathbf{R}$ ماتریس دوران با زاویهٔ $\theta$ و $\sigma_{1},\sigma_{2}$ انحراف‌معیارهای دو محور اصلی هستند. هنگامی‌ که $\sigma_{1}=\sigma_{2}$، هسته همسانگرد و در غیر این صورت ناهمسانگرد خواهد بود.

برای پوشش تنوع بیشتری از شکل‌های هسته، \lr{Real-ESRGAN} از هسته‌های \textbf{گاوسی تعمیم‌یافته} و توزیع \textbf{فلات‌گونه} نیز بهره می‌برد. تابع چگالی احتمال این دو توزیع به ترتیب عبارت‌اند از:
\[
p_{\text{gen}}(\mathbf{C}) = \frac{1}{N}\exp\!\Bigl(-\frac{1}{2}(\mathbf{C}^{\top}\boldsymbol{\Sigma}^{-1}\mathbf{C})^{\beta}\Bigr),
\quad
p_{\text{plat}}(\mathbf{C}) = \frac{1}{N}\,\frac{1}{1+(\mathbf{C}^{\top}\boldsymbol{\Sigma}^{-1}\mathbf{C})^{\beta}},
\]
که $\beta$ پارامتر شکل است. نتایج تجربی نشان داده که استفاده از این هسته‌های متنوع خروجی‌های تیزتری برای نمونه‌های واقعی تولید می‌کند، هرچند تفاوت آن‌ها با هسته‌های گاوسی ساده در اغلب نمونه‌ها اندک است.

\subsection{نویز}
دو نوع نویز رایج در نظر گرفته می‌شود:
\begin{enumerate}
	\item \textbf{نویز گاوسی افزودنی:} شدت آن با انحراف معیار ($\sigma$) توزیع گاوسی کنترل می‌شود. دو حالت نویز رنگی (نمونه‌برداری مستقل در هر کانال RGB) و نویز خاکستری (نویز یکسان برای هر سه کانال) در نظر گرفته شده است.
	\item \textbf{نویز پواسون:} برای مدل‌سازی نویز حسگر ناشی از نوسانات کوانتومی فوتون‌ها به‌کار می‌رود. شدت این نویز متناسب با شدت پیکسل تصویر است و نویز هر پیکسل مستقل از دیگری خواهد بود.
\end{enumerate}

\subsection{تغییر اندازه}
عملیات کاهش‌نمونه‌برداری عملگر پایه‌ای برای ساخت تصاویر با وضوح پایین در ابرتفکیک است. \lr{Real-ESRGAN} به‌طور عمومی‌تر هم کاهش و هم افزایش‌نمونه‌برداری (عملیات تغییر اندازه) را در نظر می‌گیرد. الگوریتم‌های متعدد تغییر اندازه شامل درونیابی ناحیه‌ای، دوخطی و دومکعبی هرکدام اثرات بصری متفاوتی ایجاد می‌کنند؛ برخی خروجی‌های تار و برخی تصاویر بیش‌ازحد تیز با آرتیفکت‌های فراجهش\LTRfootnote{Overshoot} تولید می‌نمایند. در این کار، به‌منظور ایجاد تنوع بیشتر، عملیات تغییر اندازه به‌صورت تصادفی از میان سه الگوریتم ناحیه‌ای، دوخطی و دومکعبی انتخاب می‌شود.

\subsection{فشرده‌سازی \lr{JPEG}}
فشرده‌سازی \lr{JPEG} یکی از رایج‌ترین تکنیک‌های فشرده‌سازی اتلافی است که ابتدا تصویر را به فضای رنگی \lr{YCbCr} تبدیل کرده، کانال‌های کروما را کاهش‌نمونه‌برداری می‌کند، سپس تصویر به بلوک‌های $8\times8$ تقسیم و هر بلوک با تبدیل کسینوسی گسسته دوبعدی (DCT) پردازش می‌شود. کوانتیزه‌سازی ضرایب DCT آرتیفکت‌های بلوکی ناخوشایندی ایجاد می‌کند. کیفیت فشرده‌سازی با پارامتر $q \in [0,100]$ تعیین می‌شود که مقادیر پایین‌تر نسبت فشرده‌سازی بالاتر و کیفیت پایین‌تری را نشان می‌دهد. در \lr{Real-ESRGAN} از پیاده‌سازی \lr{DiffJPEG} در \lr{PyTorch} استفاده شده است.


\subsection{مدل تخریب مرتبهٔ بالا}
\label{subsec:high_order_deg}

محدودیت اصلی مدل کلاسیک در عدم توانایی بازتولید تخریب‌های پیچیدهٔ دنیای واقعی نهفته است. تخریب‌های واقعی معمولاً حاصل ترکیب زنجیره‌ای چندین فرآیند تخریب متوالی هستند. به‌عنوان مثال، تصویری که با تلفن همراه گرفته شده ممکن است تخریب‌هایی نظیر تاری دوربین، نویز حسگر، آرتیفکت‌های تیزسازی و فشرده‌سازی JPEG را شامل شود. سپس هنگام ویرایش و بارگذاری در شبکه‌های اجتماعی، فشرده‌سازی‌های بیشتر و نویزهای غیرقابل پیش‌بینی اضافه می‌شود. این فرآیند هنگامی‌ که تصویر چندین بار در اینترنت به‌اشتراک گذاشته شود، پیچیده‌تر نیز خواهد شد.
\begin{figure}
	\centering
	\includegraphics[width=1\linewidth]{figures/ESR_HighOrder.png}
	\caption{مدل تخریب مرتبه بالا}
	\label{fig:ESR_HighOrder}
\end{figure}

این مشاهده انگیزهٔ اصلی توسعهٔ مدل تخریب مرتبهٔ بالا است. به‌جای اعمال یک‌بارهٔ مدل کلاسیک (مدل‌سازی مرتبهٔ اول)، تخریب‌ها با تکرار چندبارهٔ فرآیند تخریب کلاسیک مدل‌سازی می‌شوند:
\begin{equation}
	\label{eq:high_order}
	\mathbf{x} = \mathcal{D}_{n}(\mathbf{y}) = (\mathcal{D}_{n} \circ \cdots \circ \mathcal{D}_{2} \circ \mathcal{D}_{1})(\mathbf{y}),
\end{equation}
که هر $\mathcal{D}_{i}$ از همان ساختار مدل کلاسیک (معادلهٔ \eqref{eq:classical_deg}) پیروی می‌کند، اما با ابرپارامترهای متفاوت. از لحاظ تجربی، \lr{Real-ESRGAN} از فرآیند تخریب مرتبهٔ دوم ($n=2$) استفاده می‌کند که تعادل مناسبی میان سادگی و اثربخشی برقرار می‌سازد.

تفاوت کلیدی این رویکرد با استراتژی ترتیب‌دهی تصادفی در آن است که مدل مرتبهٔ بالا فرآیند تولید تخریب واقعی را شبیه‌سازی می‌کند و بر تکرار کامل فرآیند تخریب تأکید دارد، نه صرفاً بر جابه‌جایی تصادفی مؤلفه‌های منفرد. در رویکرد ترتیب‌دهی تصادفی، تعداد فرآیندهای تخریب ثابت است و مشخص نیست که آیا همهٔ ترتیب‌های ممکن واقعاً مفید هستند یا خیر. در مقابل، مدل‌سازی مرتبهٔ بالا انعطاف‌پذیری بیشتری دارد و مرز تخریب‌های قابل‌حل را گسترش می‌دهد.

لازم به ذکر است که عملیات کاهش‌نمونه‌برداری در معادلهٔ \eqref{eq:classical_deg} با عملیات تغییر اندازهٔ تصادفی جایگزین می‌شود تا وضوح تصویر در بازهٔ معقولی باقی بماند. نویسندگان تصریح کرده‌اند که فرآیند تخریب مرتبهٔ بالا کامل نیست و نمی‌تواند تمامی فضای تخریب دنیای واقعی را پوشش دهد؛ بلکه صرفاً مرز تخریب‌های قابل‌حل را نسبت به روش‌های پیشین گسترش می‌دهد.


\subsection{آرتیفکت‌های حلقه‌ای و فراجهش}
\label{subsec:ringing_overshoot}

آرتیفکت‌های حلقه‌ای\LTRfootnote{Ringing Artifacts} به‌صورت لبه‌های کاذب در نزدیکی گذارهای تند تصویر ظاهر می‌شوند و از لحاظ بصری شبیه نوارها یا «ارواح» در کنار لبه‌ها هستند. آرتیفکت‌های فراجهش معمولاً همراه با آرتیفکت‌های حلقه‌ای بروز می‌کنند و به‌صورت افزایش ناگهانی مقدار پیکسل در ناحیهٔ گذار لبه مشاهده می‌شوند. علت اصلی این آرتیفکت‌ها محدودیت پهنای باند سیگنال و حذف فرکانس‌های بالاست. این آرتیفکت‌ها بسیار رایج بوده و معمولاً توسط الگوریتم‌های تیزسازی، فشرده‌سازی JPEG و سایر پردازش‌ها تولید می‌شوند.
\begin{figure}
	\centering
	\includegraphics[width=1\linewidth]{figures/ESR_overshoot.png}
	\caption{آرتیفکت‌های حلقه و فراجهش در تصاویر واقعی و تولید آن با فیلتر sinc}
	\label{fig:ESR_overshoot}
\end{figure}

\lr{Real-ESRGAN} از فیلتر \lr{sinc} برای شبیه‌سازی این آرتیفکت‌ها در جفت‌داده‌های آموزشی بهره می‌برد. فیلتر \lr{sinc} یک فیلتر ایده‌آل‌سازی‌شده است که فرکانس‌های بالا را قطع می‌کند. هستهٔ فیلتر \lr{sinc} به‌صورت زیر بیان می‌شود:
\begin{equation}
	\label{eq:sinc_kernel}
	k(i,j) = \frac{\omega_{c}}{2\pi\sqrt{i^{2}+j^{2}}}\,J_{1}\!\bigl(\omega_{c}\sqrt{i^{2}+j^{2}}\bigr),
\end{equation}
که $(i,j)$ مختصات هسته، $\omega_{c}$ فرکانس قطع و $J_{1}$ تابع بسل مرتبهٔ اول از نوع اول است. با تغییر فرکانس قطع، می‌توان آرتیفکت‌های حلقه‌ای و فراجهش با شدت‌های مختلف تولید کرد.

فیلتر \lr{sinc} در دو جایگاه به‌کار گرفته می‌شود: \textbf{(۱)} در فرآیند تاری به‌عنوان جایگزینی تصادفی برای هسته‌های گاوسی و \textbf{(۲)} در آخرین مرحلهٔ ساخت داده. ترتیب آخرین فیلتر \lr{sinc} و فشرده‌سازی JPEG به‌صورت تصادفی جابه‌جا می‌شود تا فضای تخریب بزرگ‌تری پوشش داده شود. به‌عنوان مثال، برخی تصاویر ممکن است ابتدا بیش‌ازحد تیز شوند (با آرتیفکت فراجهش) و سپس فشرده‌سازی JPEG ببینند؛ در حالی‌که برخی دیگر ابتدا فشرده‌سازی JPEG و سپس تیزسازی را تجربه می‌کنند.

\begin{table}[!ht]
	\centering
	\caption{تنظیمات کلیدی فرآیند ساخت دادهٔ مصنوعی در \lr{Real-ESRGAN}. مقادیر داخل پرانتز مربوط به فرآیند تخریب مرتبهٔ دوم هستند.}
	\label{tbl:deg_settings}
	\renewcommand{\arraystretch}{1.4}
	\begin{tabular}{|l|l|}
		\hline
		\textbf{پارامتر} & \textbf{مقدار / بازه} \\ \hline
		نوع هسته‌های تاری & گاوسی (۰٫۷)، گاوسی تعمیم‌یافته (۰٫۱۵)، فلات‌گونه (۰٫۱۵) \\ \hline
		اندازهٔ هسته & تصادفی از $\{7,9,\ldots,21\}$ \\ \hline
		$\sigma$ تاری & $[0.2,\,3]$ \;($[0.2,\,1.5]$) \\ \hline
		$\beta$ شکل & $[0.5,\,4]$ (گاوسی تعمیم‌یافته)، $[1,\,2]$ (فلات‌گونه) \\ \hline
		احتمال هستهٔ \lr{sinc} & $0.1$ \\ \hline
		احتمال حذف تاری دوم & $0.2$ \\ \hline
		نوع نویز & گاوسی (۰٫۵)، پواسون (۰٫۵) \\ \hline
		$\sigma$ نویز گاوسی & $[1,\,30]$ \;($[1,\,25]$) \\ \hline
		مقیاس نویز پواسون & $[0.05,\,3]$ \;($[0.05,\,2.5]$) \\ \hline
		احتمال نویز خاکستری & $0.4$ \\ \hline
		ضریب کیفیت JPEG ($q$) & $[30,\,95]$ \\ \hline
		احتمال فیلتر \lr{sinc} نهایی & $0.8$ \\ \hline
	\end{tabular}
\end{table}

\subsection{شبکهٔ مولد}
\lr{Real-ESRGAN} از همان معماری مولد \lr{ESRGAN} بهره می‌برد، یعنی یک شبکهٔ عمیق متشکل از چندین بلوک چگال پسماند در پسماند\LTRfootnote{Residual-in-Residual Dense Block (RRDB)}. این معماری به‌دلیل استفاده از اتصالات چگال و اتصالات پسماند تودرتو، قابلیت بازنمایی بسیار بالایی دارد.

برای پشتیبانی از ضریب بزرگ‌نمایی $\times 2$ و $\times 1$ (علاوه بر $\times 4$ اصلی)، ابتدا عملیات \lr{pixel-unshuffle} (عکس \lr{pixel-shuffle}) بر روی ورودی اعمال می‌شود. این عملیات اندازهٔ فضایی تصویر ورودی را کاهش داده و اطلاعات را در بعد کانال بازآرایی می‌کند. در نتیجه، بخش اعظم محاسبات در فضای وضوح کمتر انجام می‌شود که مصرف حافظهٔ GPU و هزینهٔ محاسباتی را به‌طور قابل‌ملاحظه‌ای کاهش می‌دهد.
\begin{figure}
	\centering
	\includegraphics[width=1\linewidth]{figures/ESRGAN_Arch.png}
	\caption{معماری شبکه مولد ESRGAN}
	\label{fig:ESR_ESRGAN_Arch}
\end{figure}

\subsection{تمایزگر \lr{U-Net} با نرمال‌سازی طیفی}
از آنجا که \lr{Real-ESRGAN} فضای تخریب بسیار بزرگ‌تری نسبت به \lr{ESRGAN} را هدف قرار می‌دهد، طراحی اصلی تمایزگر در \lr{ESRGAN} (تمایزگر سبک \lr{VGG}) دیگر کافی نیست. تمایزگر در \lr{Real-ESRGAN} نیازمند قدرت تمایز بیشتری برای خروجی‌های پیچیدهٔ آموزشی است. علاوه بر تمایز سبک‌های سراسری، تمایزگر باید بازخورد گرادیانی دقیقی برای بافت‌های محلی فراهم کند.

با الهام از \cite{schonfeld2020unet}, تمایزگر سبک \lr{VGG} در \lr{ESRGAN} به یک طراحی \lr{U-Net} با اتصالات میان‌بُری ارتقا یافته است. خروجی \lr{U-Net} مقادیر واقعی‌بودن\LTRfootnote{Realness Values} را برای هر پیکسل تولید می‌کند و می‌تواند بازخورد جزئی پیکسل‌به‌پیکسل را به مولد ارائه دهد.

ساختار \lr{U-Net} و تخریب‌های پیچیده، ناپایداری آموزش را افزایش می‌دهند. برای رفع این مشکل از نرمال‌سازی طیفی\LTRfootnote{Spectral Normalization} \cite{ZhangDenoise2017} استفاده شده است. مشاهده شده که نرمال‌سازی طیفی علاوه بر پایدارسازی دینامیک آموزش، به کاهش آرتیفکت‌های بیش‌تیز و ناخوشایند ناشی از آموزش GAN نیز کمک می‌کند. با این تنظیمات، تعادل مطلوبی میان بهبود جزئیات محلی و سرکوب آرتیفکت‌ها حاصل می‌شود.

\begin{figure}
	\centering
	\includegraphics[width=1\linewidth]{figures/ESR_Unet.png}
	\caption{تمایزگر U-Net}
	\label{fig:ESR_ESR_Unet}
\end{figure}

\subsection{فرآیند آموزش}
\label{subsec:training}

آموزش \lr{Real-ESRGAN} در دو مرحله انجام می‌شود:

\noindent\textbf{مرحلهٔ اول (آموزش \lr{Real-ESRNet}):}
ابتدا یک مدل معطوف به PSNR با تابع هزینهٔ $L_1$ آموزش داده می‌شود. این مدل \lr{Real-ESRNet} نامیده شده و از وزن‌های پیش‌آموزش‌دیدهٔ \lr{ESRGAN} به‌عنوان مقداردهی اولیه بهره می‌برد. آموزش به‌مدت $1{,}000$K تکرار با نرخ یادگیری $2 \times 10^{-4}$ انجام می‌شود.

\noindent\textbf{مرحلهٔ دوم (آموزش \lr{Real-ESRGAN}):}
مدل \lr{Real-ESRNet} به‌عنوان مقداردهی اولیهٔ مولد استفاده شده و سپس \lr{Real-ESRGAN} با ترکیبی از سه تابع هزینه آموزش داده می‌شود:
\begin{equation}
	\label{eq:total_loss}
	\mathcal{L}_{\text{total}} = \underbrace{\mathcal{L}_{L_1}}_{\text{وزن } 1} + \underbrace{\mathcal{L}_{\text{perceptual}}}_{\text{وزن } 1} + \underbrace{\mathcal{L}_{\text{GAN}}}_{\text{وزن } 0.1},
\end{equation}
که $\mathcal{L}_{L_1}$ هزینهٔ بازسازی پیکسلی، $\mathcal{L}_{\text{perceptual}}$ هزینهٔ ادراکی و $\mathcal{L}_{\text{GAN}}$ هزینهٔ رقابتی است. هزینهٔ ادراکی از نقشه‌های ویژگی $\{\text{conv}1,\ldots,\text{conv}5\}$ (قبل از فعال‌سازی) شبکهٔ \lr{VGG19} پیش‌آموزش‌دیده با وزن‌های $\{0.1, 0.1, 1, 1, 1\}$ محاسبه می‌شود. آموزش مرحلهٔ دوم به‌مدت $400$K تکرار با نرخ یادگیری $1 \times 10^{-4}$ صورت می‌گیرد.

\begin{table}[!ht]
	\centering
	\caption{جزئیات آموزش \lr{Real-ESRGAN}.}
	\label{tbl:training_details}
	\renewcommand{\arraystretch}{1.4}
	\begin{tabular}{|l|l|}
		\hline
		\textbf{پارامتر} & \textbf{مقدار} \\ \hline
		مجموعه‌دادهٔ آموزشی & \lr{DIV2K + Flickr2K + OST} \\ \hline
		اندازهٔ پچ HR & $256 \times 256$ \\ \hline
		اندازهٔ دسته & $48$ (بر $4$ GPU) \\ \hline
		بهینه‌ساز & \lr{Adam} \\ \hline
		تکرارهای \lr{Real-ESRNet} & $1{,}000$K \\ \hline
		نرخ یادگیری \lr{Real-ESRNet} & $2 \times 10^{-4}$ \\ \hline
		تکرارهای \lr{Real-ESRGAN} & $400$K \\ \hline
		نرخ یادگیری \lr{Real-ESRGAN} & $1 \times 10^{-4}$ \\ \hline
		میانگین‌گیری متحرک نمایی (EMA) & فعال \\ \hline
	\end{tabular}
\end{table}

\subsection{استخر جفت‌دادهٔ آموزشی}
تمامی فرآیندهای تخریب در \lr{PyTorch} با شتاب‌دهی CUDA پیاده‌سازی شده‌اند تا جفت‌داده‌های آموزشی به‌صورت برخط\LTRfootnote{On-the-fly} ساخته شوند. با این حال، پردازش دسته‌ای تنوع تخریب‌های مصنوعی درون یک دسته را محدود می‌کند. برای مثال، نمونه‌های یک دسته نمی‌توانند ضرایب تغییر اندازهٔ متفاوتی داشته باشند. برای رفع این محدودیت، از یک استخر جفت‌دادهٔ آموزشی با اندازهٔ $180$ استفاده شده است. در هر تکرار، نمونه‌های آموزشی به‌صورت تصادفی از این استخر انتخاب می‌شوند تا دستهٔ آموزشی تشکیل شود.

\subsection{تیزسازی تصاویر حقیقی در حین آموزش}
یک ترفند آموزشی مؤثر برای بهبود بصری تیزی خروجی، تیزسازی تصاویر حقیقی (زمینه‌ای) در حین آموزش است. روش متداول تیزسازی پس‌پردازش مانند ماسک‌گذاری ناتیز\LTRfootnote{Unsharp Masking (USM)} تمایل به ایجاد آرتیفکت‌های فراجهش دارد. تجربیات نشان داده که تیزسازی تصاویر حقیقی در مرحلهٔ آموزش، تعادل بهتری میان تیزی و سرکوب آرتیفکت‌های فراجهش ایجاد می‌کند. مدل آموزش‌دیده با تصاویر حقیقی تیزشده با نام \lr{Real-ESRGAN+} معرفی شده است.


\subsection{مقایسه با روش‌های پیشین}
\label{subsec:comparisons}

\lr{Real-ESRGAN} با چندین روش پیشرفتهٔ ابرتفکیک شامل \lr{ESRGAN} \cite{wang2018esrganenhancedsuperresolutiongenerative}، \lr{DAN} \cite{luo2020dan}، \lr{CDC} \cite{wei2020cdc}، \lr{RealSR} \cite{ji2020realsr} و \lr{BSRGAN} \cite{zhang2021bsrgan} مقایسه شده است. ارزیابی بر روی مجموعه‌دادهٔ‌های متنوع تصاویر دنیای واقعی شامل \lr{RealSR}، \lr{DRealSR}، \lr{OST300}، \lr{DPED}، \lr{ADE20K} و تصاویر اینترنتی انجام شده است.

از آنجا که معیارهای موجود کیفیت ادراکی نمی‌توانند ترجیحات واقعی ادراک بشری را در مقیاس ریزدانه به‌خوبی منعکس کنند، نویسندگان بر ارزیابی بصری تأکید ویژه‌ای کرده‌اند. با این حال، نتایج کمّی معیار NIQE \cite{mittal2013niqe} نیز گزارش شده است.

\begin{table}[!ht]
	\centering
	\caption{امتیازات NIQE بر روی مجموعه‌دادهٔ‌های مختلف (مقادیر کمتر بهتر هستند). بهترین نتیجه \textbf{پررنگ} شده است.}
	\label{tbl:niqe_comparison}
	\renewcommand{\arraystretch}{1.3}
	\resizebox{\textwidth}{!}{%
	\begin{tabular}{|l|c|c|c|c|c|c|c|c|}
		\hline
		\textbf{مجموعه‌داده} & \lr{Bicubic} & \lr{ESRGAN} & \lr{DAN} & \lr{RealSR} & \lr{CDC} & \lr{BSRGAN} & \lr{Real-ESRGAN} & \lr{Real-ESRGAN+} \\ \hline
		\lr{RealSR-Canon} & 6.13 & 6.77 & 6.53 & 6.87 & 6.15 & 5.75 & 4.59 & \textbf{4.53} \\ \hline
		\lr{RealSR-Nikon} & 6.36 & 6.75 & 6.61 & 6.74 & 6.33 & 5.99 & 5.08 & \textbf{5.02} \\ \hline
		\lr{DRealSR}      & 6.58 & 8.63 & 7.07 & 7.72 & 6.64 & 6.14 & 4.98 & \textbf{4.85} \\ \hline
		\lr{DPED-iPhone}  & 6.01 & 5.74 & 6.14 & 5.59 & 6.27 & 5.99 & 5.44 & \textbf{5.26} \\ \hline
		\lr{OST300}       & 4.44 & 3.52 & 5.02 & 4.57 & 4.74 & 4.17 & 2.87 & \textbf{2.82} \\ \hline
	\end{tabular}%
	}
\end{table}

نتایج بصری نشان می‌دهد که \lr{Real-ESRGAN} از لحاظ حذف آرتیفکت‌ها و بازسازی جزئیات بافت، عملکرد بهتری نسبت به روش‌های پیشین دارد. به‌طور مشخص:
\begin{itemize}
	\item \textbf{حذف آرتیفکت‌های فراجهش:} نمونه‌هایی که حاوی آرتیفکت‌های فراجهش (لبه‌های سفید اطراف حروف) هستند، توسط روش‌هایی مانند \lr{DAN} و \lr{BSRGAN} به‌صورت تقویت‌شده باقی می‌مانند؛ در حالی‌که \lr{Real-ESRGAN} با شبیه‌سازی این آرتیفکت‌ها در آموزش، آن‌ها را به‌طور مؤثری حذف می‌کند.
	\item \textbf{حذف تخریب‌های پیچیدهٔ ناشناخته:} اغلب الگوریتم‌ها قادر به حذف مؤثر تخریب‌های پیچیده و ناشناخته نیستند، اما \lr{Real-ESRGAN} به‌واسطهٔ آموزش با فرآیند تخریب مرتبهٔ دوم، این قابلیت را دارد.
	\item \textbf{بازسازی بافت‌های طبیعی:} \lr{Real-ESRGAN} بافت‌های واقع‌گرایانه‌تری (مانند بافت آجر، کوه و درخت) برای نمونه‌های واقعی بازسازی می‌کند؛ در حالی‌که سایر روش‌ها یا در حذف تخریب‌ها ناکام می‌مانند یا بافت‌های غیرطبیعی اضافه می‌کنند.
\end{itemize}


\subsection{مطالعات تفکیکی}
\label{subsec:ablation}

نویسندگان مطالعات تفکیکی\LTRfootnote{Ablation Studies} جامعی برای بررسی اثر هر مؤلفه ارائه کرده‌اند:

\subsection{اثر مدل تخریب مرتبهٔ دوم}
مدل‌هایی که با مدل تخریب کلاسیک (مرتبهٔ اول) آموزش دیده‌اند، قادر به حذف مؤثر نویز یا تاری در نمونه‌هایی مانند بافت دیوار یا مزرعهٔ گندم نیستند. در مقابل، \lr{Real-ESRNet} آموزش‌دیده با تخریب مرتبهٔ دوم، این موارد را با موفقیت مدیریت می‌کند. این نتیجه نشان‌دهندهٔ اهمیت حیاتی مدل‌سازی تخریب مرتبهٔ بالا در مواجهه با تخریب‌های دنیای واقعی است.

\subsection{اثر فیلتر \lr{sinc}}
بدون استفاده از فیلتر \lr{sinc} در آموزش، نتایج بازسازی‌شده آرتیفکت‌های حلقه‌ای و فراجهش موجود در تصاویر ورودی را تقویت می‌کنند، به‌ویژه اطراف متن و خطوط. در مقابل، مدل‌هایی که با فیلتر \lr{sinc} آموزش دیده‌اند، قادر به حذف این آرتیفکت‌ها هستند.
\subsection{اثر تمایزگر \lr{U-Net} با نرمال‌سازی طیفی}
ارزیابی سه پیکربندی مختلف تمایزگر نتایج زیر را نشان می‌دهد:
\begin{enumerate}
	\item \textbf{تنظیمات \lr{ESRGAN} (تمایزگر سبک \lr{VGG}):} ناتوانی در بازسازی بافت‌های ریز (آجر و بوته) و ایجاد آرتیفکت‌های ناخوشایند.
	\item \textbf{تمایزگر \lr{U-Net} بدون نرمال‌سازی طیفی:} بهبود جزئیات محلی، اما ایجاد بافت‌های غیرطبیعی و افزایش ناپایداری آموزش.
	\item \textbf{تمایزگر \lr{U-Net} با نرمال‌سازی طیفی:} بهبود بافت‌های بازسازی‌شده همراه با پایدارسازی دینامیک آموزش.
\end{enumerate}

\subsection{اثر هسته‌های تاری پیچیده‌تر}
حذف هسته‌های گاوسی تعمیم‌یافته و فلات‌گونه از فرآیند ساخت داده نشان داد که در برخی نمونه‌های واقعی، مدل نمی‌تواند تاری را حذف کرده و لبه‌های تیز بازیابی کند. با این حال، در اکثر نمونه‌ها تفاوت‌ها اندک بوده و هسته‌های گاوسی با فرآیند تخریب مرتبهٔ بالا به‌تنهایی فضای وسیعی از تاری‌های واقعی را پوشش می‌دهد. نویسندگان با مشاهدهٔ بهبود اندک اما قابل‌مشاهده، استفاده از هسته‌های پیچیده‌تر را حفظ کرده‌اند.


\subsection{محدودیت‌ها}
\label{subsec:limitations}

با وجود توانایی بالای \lr{Real-ESRGAN} در بازسازی اکثر تصاویر دنیای واقعی، این روش محدودیت‌هایی نیز دارد:
\begin{enumerate}
	\item \textbf{خطوط پیچ‌خورده:} برخی تصاویر بازسازی‌شده (به‌ویژه صحنه‌های ساختمانی و داخلی) به‌دلیل مشکلات نام‌مستعار\LTRfootnote{Aliasing} خطوط پیچ‌خورده‌ای نشان می‌دهند.
	\item \textbf{آرتیفکت‌های آموزش GAN:} آموزش GAN در برخی نمونه‌ها آرتیفکت‌های ناخوشایندی ایجاد می‌کند.
	\item \textbf{تخریب‌های خارج از توزیع:} مدل قادر به حذف تخریب‌های پیچیدهٔ خارج از توزیع آموزشی نیست و حتی ممکن است این آرتیفکت‌ها را تقویت کند.
\end{enumerate}
نویسندگان تأکید کرده‌اند که این محدودیت‌ها تأثیر قابل‌توجهی بر کاربرد عملی \lr{Real-ESRGAN} دارند و رفع آن‌ها نیازمند پژوهش‌های آتی است.


\subsection{نتیجه‌گیری و جمع‌بندی}
\label{subsec:realesrgan_summary}

\lr{Real-ESRGAN} با سه نوآوری اصلی، مرزهای ابرتفکیک کور دنیای واقعی را گسترش داده است:
\begin{itemize}
	\item \textbf{مدل تخریب مرتبهٔ بالا:} فرآیند تخریب مرتبهٔ دوم به‌همراه فیلترهای \lr{sinc} برای شبیه‌سازی آرتیفکت‌های حلقه‌ای و فراجهش، فضای تخریب‌های قابل‌حل را نسبت به مدل کلاسیک به‌طور قابل‌ملاحظه‌ای گسترش می‌دهد.
	\item \textbf{تمایزگر \lr{U-Net} با نرمال‌سازی طیفی:} قدرت تمایز بالاتر و بازخورد پیکسل‌به‌پیکسل، همراه با پایدارسازی دینامیک آموزش.
	\item \textbf{آموزش با دادهٔ کاملاً مصنوعی:} مدل بدون نیاز به دادهٔ جفت‌شدهٔ واقعی، تنها با جفت‌داده‌های مصنوعی آموزش می‌بیند و با این حال قادر به بازسازی اکثر تصاویر دنیای واقعی است.
\end{itemize}

نتایج تجربی نشان‌دهندهٔ عملکرد بصری برتر \lr{Real-ESRGAN} نسبت به روش‌های پیشین بر مجموعه‌دادهٔ‌های متنوع واقعی است. پیاده‌سازی کارآمد فرآیند ساخت دادهٔ مصنوعی به‌صورت برخط و تکنیک‌هایی مانند استخر جفت‌داده و تیزسازی تصاویر حقیقی، این مدل را به یک ابزار عملی برای ابرتفکیک تصاویر دنیای واقعی تبدیل کرده است. با این حال، محدودیت‌هایی مانند خطوط پیچ‌خورده، آرتیفکت‌های GAN و ناتوانی در مواجهه با تخریب‌های خارج از توزیع، فضای پژوهشی مهمی برای کارهای آینده باقی گذاشته‌اند.
